\section*{Заключение}
\addcontentsline{toc}{section}{Заключение}

В результате выполнения курсового проекта была спроектирована и успешно внедрена веб-платформа маркетплейса электронной техники. Практическая 
реализация подтвердила актуальность создания торговых площадок, предоставляющих представителям малого и среднего бизнеса специализированные инструменты 
для дистанционной торговли, а розничным покупателям — удобную систему поиска и агрегации цифровых товаров.

На основе полученных при разработке данных можно сформулировать \textbf{выводы}.

\begin{enumerate}
    \item Использование раздельной архитектуры клиент-серверного взаимодействия обеспечило высокую производительность приложения и гарантировало безопасность бизнес-логики.
    \item Проектирование инфологической модели с применением СУБД PostgreSQL и типа данных JSONB доказало свою высокую эффективность - это позволило организовать 
	надежное хранение неструктурированных спецификаций товаров и реализовать алгоритмы динамической клиентской фильтрации без избыточного усложнения схемы БД.
    \item Внедрение ролевой модели доступа на основе криптографически подписанных токенов дало возможность безопасно объединить в единой экосистеме витрину товаров для покупателей 
	и изолированные панели управления для продавцов и администраторов с функциями генерации финансовой отчетности.
\end{enumerate}

\textbf{Рекомендации и предложения.} Для успешного внедрения разработанной платформы в реальную коммерческую среду рекомендуется реализовать API-интеграцию с внешними учетными 
системами по типу «1С:Предприятие» для автоматизации документооборота продавцов. Также целесообразно подключить шлюзы сторонних логистических операторов (например, СДЭК) для 
автоматического расчета стоимости доставки и отслеживания отправлений в личном кабинете покупателя в режиме реального времени.

\textbf{Перспективы углубления темы.} В дальнейшем исследовании планируется развитие функциональности системы за счет интеграции 
эквайринговых шлюзов для проведения защищенных онлайн-транзакций, разработки микросервиса рекомендательных алгоритмов на основе поведенческого анализа покупателей, а также 
расширения аналитического модуля для продавцов с прогнозируемым спросом на продукт.

При разработке пояснительной записки к курсовому проекту строго соблюдались требования~\cite{gost_report}, определяющие структуру и правила оформления 
научно-исследовательских работ.

